\section{Введение}

Контейнеризация стала стандартом поставки серверного программного обеспечения: образы дают воспроизводимую среду, а декларативные манифесты снижают дрейф между стендами. Вместе с тем инженеру, который ведёт несколько сервисов на собственных или арендованных машинах, не всегда нужен полный масштаб промышленного оркестратора. Именно из практической потребности «управлять инфраструктурой так же предсказуемо, как кодом приложения», но без избыточной сложности, и возник проект \textbf{Chronos}.

\subsection{Практическая мотивация}

В экосистеме Docker давно существуют зрелые инструменты. \textbf{Kubernetes} даёт мощную модель деплоя, сетевых политик, масштабирования и самовосстановления, однако требует отдельного кластера, специалистов по эксплуатации и существенных накладных расходов на обучение и сопровождение. Для небольшой команды или личного стенда это часто избыточно: стоимость владения оказывается выше, чем выигрыш от функций, которые фактически не используются.

\textbf{Portainer} и схожие панели закрывают другой класс задач: наглядное управление контейнерами и compose-стеками через веб-интерфейс. При этом конфигурация по-прежнему в основном живёт в YAML-файлах и ручных действиях оператора; связка «репозиторий приложения --- декларативное описание инфраструктуры --- CI/CD» не всегда выражена единым программным API, которым удобно пользоваться из кода сборки или интеграционных тестов.

Опыт облачных платформ, в частности подходов, близких к экосистеме \textbf{Microsoft Azure} (инфраструктура и параметры приложения как управляемый из кода артефакт с проверкой на этапе сборки), показал привлекательность модели: разработчик описывает среду рядом с логикой сервиса, версионирует изменения тем же процессом, что и код, и может прогонять проверки до выката. Перенести эту идею в мир Docker без обязательного перехода на Kubernetes естественно через \textbf{Docker Compose}: формат широко распространён, хорошо документирован и поддерживается локально и в CI.

\subsection{Идея Chronos}

Платформа Chronos объединяет три уровня:

\begin{itemize}[leftmargin=1.5cm]
  \item \textbf{Chronos.Core} --- библиотека на .NET с Fluent API для описания compose-проектов, генерации YAML, валидации, локального и удалённого запуска; поддержка кодовых тестов и job-ов (\texttt{[Test]}, \texttt{[Job]}), прикрепляемых к сервису через \texttt{UseTests}/\texttt{UseJobs};
  \item \textbf{Chronos.Agent} --- сервис на машине с Docker Engine, который принимает манифесты, выполняет \texttt{docker compose}, отдаёт статус, логи и снимки томов;
  \item \textbf{Chronos.Master} (и веб-клиент \textbf{Chronos.Master.Web}) --- центральный узел: реестр агентов, аудит, проксирование запросов UI к агентам, политики резервного копирования томов в объектное хранилище, динамическая генерация TCP-конфигурации для HAProxy.
\end{itemize}

Таким образом, Chronos не заменяет Docker, а надстраивает над ним управляющий контур с сохранением compose-парадигмы.

\subsection{Цель и задачи работы}

\textbf{Цель работы} --- разработать и описать программный комплекс Chronos для централизованного управления Docker Compose-проектами на нескольких узлах с поддержкой наблюдаемости, резервного копирования томов и динамической TCP-маршрутизации.

\textbf{Задачи:}
\begin{enumerate}[leftmargin=1.5cm]
  \item Проанализировать ограничения существующих классов решений (Kubernetes, панели управления, «голый'' Compose) для целевого сценария.
  \item Сформулировать требования к модульной архитектуре (Core, Agent, Master) и к нефункциональным свойствам (наблюдаемость, воспроизводимость).
  \item Спроектировать взаимодействие компонентов: публикация проектов, heartbeat агентов, политики бэкапа томов, реестр TCP-маршрутов HAProxy.
  \item Реализовать серверные и клиентские части, развернуть контрольный стенд в Docker.
  \item Провести функциональную проверку ключевых сценариев и зафиксировать ограничения MVP.
\end{enumerate}

\subsection{Объект, предмет и методы}

\textbf{Объект исследования} --- распределённая программная система управления контейнерными приложениями на базе Docker Compose.

\textbf{Предмет исследования} --- методы построения агентной архитектуры, алгоритмы согласования состояния маршрутов и политик резервного копирования с объектным хранилищем S3-совместимого типа, а также интеграция с HAProxy и стеком Loki/Grafana.

\textbf{Методы:} системный анализ, объектно-ориентированное проектирование, экспериментальное развёртывание на локальном стенде, анализ логов и конфигураций.

\subsection{Научная новизна и практическая значимость}

К практической новизне относятся: связка «Fluent API $\leftrightarrow$ YAML $\leftrightarrow$ удалённый агент'' с возможностью обратного разбора compose-файла в цепочку вызовов билдера; централизованные политики снимков томов с учётом свободного места на узле агента и ретенции копий в бакете; динамическое обновление фрагмента конфигурации HAProxy с учётом эксплуатационных деталей (кодировка UTF-8 без BOM, перезагрузка прокси).

\textbf{Практическая значимость} --- готовый к использованию контур для локальной и небольшой командной инфраструктуры, снижающий долю ручных операций при публикации сервисов и при публикации TCP-портов наружу.

\subsection{Структура работы}

Работа состоит из введения, глав с анализом предметной области, архитектуры и проектирования, реализации, результатов испытаний, заключения, приложений и списка использованных источников. В приложениях приведены вспомогательные таблицы сценариев и справочные материалы.
